% ENG 2440 Assignment 3 report. Build: tectonic main.tex   (figures are read from ../figures/)
\documentclass[11pt,a4paper]{article}
\usepackage[margin=2cm]{geometry}            % \textwidth = 483.697 pt = 6.693 in; figures are drawn at
\usepackage{newtxtext,newtxmath}             % that width and never rescaled. Times, as in the figures.
\usepackage{graphicx,booktabs,multirow,array,xcolor,microtype}
\usepackage[font=footnotesize,labelfont=bf,skip=3pt]{caption}
\usepackage[compact]{titlesec}
\usepackage{enumitem,placeins,capt-of}
\usepackage[hidelinks]{hyperref}

\graphicspath{{../figures/}}
\newcommand{\fig}[2][1]{\includegraphics[width=#1\dimexpr 6.693in\relax]{#2}}
\newcommand{\repourl}{https://github.com/DannyZ061205/ENG2440-A3}
\setlist{nosep,leftmargin=*}
\setlength{\emergencystretch}{1.5em}
\setcounter{topnumber}{3}\setcounter{bottomnumber}{2}\setcounter{totalnumber}{5}
\renewcommand{\topfraction}{0.92}\renewcommand{\bottomfraction}{0.65}\renewcommand{\textfraction}{0.06}
\setlength{\textfloatsep}{6pt plus 2pt minus 2pt}
\setlength{\floatsep}{6pt plus 2pt minus 2pt}
\setlength{\intextsep}{6pt plus 2pt minus 2pt}
\titlespacing*{\section}{0pt}{6pt}{2pt}
\titlespacing*{\subsection}{0pt}{6pt}{2pt}
\setlength{\parskip}{1pt plus 1pt}

\begin{document}

\begin{center}
{\Large\bfseries Predicting blood--brain barrier penetration from molecular graphs}\\[3pt]
Zichen Zhao (25011189) \quad·\quad ENG 2440 Medical Imaging \& AI in Healthcare, Assignment 3\\[1pt]
Code and executed notebook: \url{\repourl}
\end{center}
\suppressfloats[t]  % page 1 only: a top float here would be typeset above the title block

\noindent\textbf{Summary.} A molecule is a graph, so this report builds one for every compound in the
BBBP dataset \cite{wu2018,martins2012} and asks whether training a graph neural network on that
representation helps predict the blood--brain barrier label. Cleaning removes 11 unparseable SMILES,
collapses duplicates by canonical SMILES and deletes 10 molecules carrying contradictory labels, leaving
1,965 molecules, 76.3\% positive. A bond-aware graph neural network (3 GINEConv layers, 242,564
parameters) is compared against Morgan fingerprints with logistic regression under a stratified random
split and a Bemis--Murcko scaffold split \cite{bemis1996}. The headline result is negative: \emph{the
graph network does not beat the fingerprint baseline}. On the untouched test sets the baseline reaches
ROC-AUC 0.919 against the network's $0.903 \pm 0.009$ (random) and 0.876 against $0.866 \pm 0.011$
(scaffold), gaps smaller than the baseline's own 95\% bootstrap interval ($\pm 0.04$). What does move the
score is the split: both models drop by 0.037--0.043 when analogs are separated, three to four times
the seed-to-seed spread. The label is a curated literature annotation, so every number measures agreement
with that annotation, not biological fact.

\section{Task and data}

\textbf{Task.} The input is one compound written as a SMILES string, a line of text spelling out which
atoms a molecule contains and how they are bonded. That string becomes a molecular graph in which each
heavy atom (every atom except hydrogen) is a node carrying nine categorical features and each bond an
edge carrying three. The output is one probability per molecule, which a single threshold turns into a
predicted class. The unit of analysis is the molecule: every metric counts molecules and every split
divides them. The intended use is to rank candidates so laboratory effort goes to the more promising
ones.

\textbf{What the label means.} \texttt{p\_np}~$=1$ marks a compound that curated literature records as
penetrating the barrier, 0 one recorded as non-penetrating \cite{martins2012}. It is an annotation
assembled from published studies differing in method, species and dose, not a measurement made here and
not a statement about any patient (Section~\ref{sec:clean}).

\textbf{What the molecules look like.} Figure~\ref{fig:molecules} shows six compounds chosen by a fixed
rule rather than by eye. Size varies by more than an order of magnitude---2 to 132 atoms, median 23, and a
median of 25 bonds (Figure~\ref{fig:sizes})---which is why the model pools atom vectors into a
fixed-length molecule vector instead of assuming a fixed atom count.

\begin{figure}[tb]
\centering\fig{A_molecules}
\caption{\textbf{Both classes contain drug-like molecules of similar size, so the label is not a proxy for
molecular weight.} Six compounds chosen by a fixed rule: within each class, the molecules at the 25th,
50th and 75th percentile of atom count. Top row penetrating, bottom row non-penetrating. Structures are
black on white with hydrogens implicit. The tag gives the true class by shape, filled for penetrating
and outlined for non-penetrating, leaving red and petrol to mean the two models. Columns are the three
percentiles; each panel gives the full canonical SMILES and the atom and bond counts.}
\label{fig:molecules}
\end{figure}

\section{Cleaning the dataset before any split}
\label{sec:clean}

The distributed CSV holds 2,050 records with no missing or non-binary labels. RDKit \cite{rdkit} rejects
11 SMILES strings, which is why the PyTorch Geometric loader \cite{fey2019} returns 2,039 graphs, none of
them empty. Canonicalizing exposes what raw strings hide: 1,975 distinct molecules, so 60 compounds
appear more than once under 64 duplicate records (\texttt{n(ccc1)c(c1)CCNC} and \texttt{c1cc(CCNC)ncc1}
are one molecule).

Ten of those canonical molecules carry \emph{both} labels in different records. All 20 records in those
groups are deleted, because there is no way to decide which label is right and keeping either would plant
a known-wrong label in training or, worse, in a test set. Collapsing the consistent duplicates and
removing the contradictory groups leaves \textbf{1,965 molecules, 1,500 of them positive (0.763)}, and
both splits are built only after this step. Those 10 contradictions are the clearest available measure of
how far the label itself can be trusted (Section~\ref{sec:limits}).

Hydrates and salts are kept as written, so a few graphs have more than one connected component and
their atom counts include those fragments. Messages never cross between components.

\begin{figure}[tb]
\centering\fig{A_sizes}
\caption{\textbf{Molecule size is right-skewed but narrow, and the two classes occupy the same size range.}
Atoms (a) and bonds (b) per molecule, and heavy-atom count by class (c), $n = 1{,}965$. Gray dashed lines
mark the medians. Axes are clipped at the 99.5th percentile, leaving 10 molecules beyond each
(largest: 132 atoms). Panel (c) plots each class as a share of itself, since the classes are 76:24;
penetrating is filled, non-penetrating outlined, as in Figure~\ref{fig:molecules}. Nothing here is a
model, so nothing carries a model color: red and petrol are reserved for those.}
\label{fig:sizes}
\end{figure}

\section{Graphs, features and batching}

\textbf{What the graph stores.} Each molecule becomes a table of atoms, a list of which atoms are
bonded, and a table describing those bonds (\texttt{x}, \texttt{edge\_index}, \texttt{edge\_attr});
Figure~\ref{fig:views} shows one compound in both views. Every atom contributes nine descriptors and
every bond three, and each bond is stored twice, once per direction, so a message can travel either way.
All twelve columns are \emph{category codes}, not quantities: code 7 in the element column means
nitrogen, not seven of something, so arithmetic on these integers is meaningless.
Table~\ref{tab:features} names them.

\textbf{How the codes are encoded.} Each column gets its own trainable embedding table and an atom's
nine looked-up vectors are summed into one 128-dimensional vector. Embeddings beat one-hot here for two
reasons: the element column alone has 119 categories, so one-hot would give a wide, nearly empty input,
and an embedding can learn that chlorine and bromine behave alike, whereas one-hot holds every category
equidistant. Either way the codes stay categorical and the raw integer never acts as a number.

\textbf{Batching.} Molecules differ in size, so they cannot be stacked like fixed-size images. PyTorch
Geometric joins a batch into one disconnected graph and records each atom's molecule in a \texttt{batch}
vector; pooling averages within molecules, so one prediction still comes out per molecule. A batch of 128
gives \texttt{x} of shape $(2957, 9)$ and \texttt{edge\_index} of shape $(2, 6358)$, and assertions
confirm that all 1,965 molecules and 47,599 atoms survive batching.

\begin{figure}[tb]
\centering\fig{B_views}
\caption{\textbf{The two views hold the same information; the graph view is simply explicit about it.}
Ornidazole as (a) a chemical structure and (b) the node--edge graph the model reads, on the same 2D
coordinates so each atom sits in the same place in both. The drawing hides carbons at line junctions and
hydrogens entirely; the graph draws every heavy atom as a node and keeps hydrogens only as the
\texttt{num\_hs} feature. Line weight encodes bond order and aromatic bonds are dashed, as the key in (b)
states; the model receives these as category codes, not geometry.}
\label{fig:views}
\end{figure}

\section{Two splits, and why the choice matters}

Both strategies divide the same 1,965 molecules 70:15:15 with one fixed seed (0). The \textbf{stratified
random split} shuffles each class separately, holding the positive proportion at 0.763 in all three
subsets. The \textbf{scaffold split} reduces each molecule to its Bemis--Murcko scaffold, its ring systems
plus linkers with side chains stripped \cite{bemis1996}, then assigns whole scaffold groups largest-first,
so molecules sharing a scaffold stay together and the held-out subsets are built from scaffolds the model
never saw. The 95 molecules with no ring system share the empty scaffold and are treated as one group,
preserving the same guarantee for acyclic molecules; no cleaned molecule has an uncomputable scaffold.

Both splits cover every molecule once, have pairwise-disjoint indices and share no canonical SMILES
across subsets; the scaffold split also has disjoint scaffold sets (430/295/295), and counts are in
Table~\ref{tab:splits}.

\textbf{Why a random split flatters a model.} A random split decides where a molecule goes without
looking at what it is, so a drug and its close analog can land one in training and one in test, and
predicting the test molecule then needs no generalization. This is measurable rather than rhetorical.
Figure~\ref{fig:splits}(b) gives each test molecule's Tanimoto similarity to its most similar training
molecule, Tanimoto being the fraction of bits set in both fingerprints out of those set in either. Under
the random split 43.9\% of test molecules have a training neighbor above Tanimoto 0.6 (median 0.566);
under the scaffold split 21.7\% do (median 0.410). Although the scaffold split is not stratified, all
three subsets stay near 0.76 positive, so the extra difficulty comes from chemistry, not a shifted prior.

Scaffold coverage is concentrated (Figure~\ref{fig:splits}(a)): of 1,020 distinct scaffolds, 777 hold
one molecule, while the 10 largest cover 22\% of the dataset and benzene alone covers 128. Filling
largest-first therefore leaves validation and test made entirely of single-molecule scaffolds.

\begin{figure}[tb]
\centering\fig{C_splits}
\caption{\textbf{The scaffold split removes the close training analogs that a random split leaves
behind.} (a) Cumulative share of the 1,965 molecules covered by the top $k$ scaffolds, largest first
(log axis), with the 10th and 100th largest scaffolds marked. (b) Cumulative distribution of each test
molecule's Morgan-fingerprint Tanimoto similarity to its nearest training molecule; the scaffold curve
lies entirely to the left, so its test molecules sit further from anything the model saw. The dashed
guide marks Tanimoto 0.6. Split is carried by lightness, the scaffold split dark and the random split
light. $n = 294$ (random), $n = 295$ (scaffold).}
\label{fig:splits}
\end{figure}

\section{The two models}

\textbf{Fingerprint baseline.} Morgan/ECFP fingerprints \cite{rogers2010} with \textbf{radius 2} and
\textbf{2,048 bits}: each atom's neighborhood out to two bonds is hashed and that bit set, so bit $k$ on
means the molecule contains a fragment hashing to $k$. The vector records which fragments are present but
not how many, where, or how they connect, and distinct fragments can collide in one bit; molecules set
43.0 bits on average. Logistic regression is fitted on the training subset only with
\texttt{class\_weight="balanced"}, without which the classifier is rewarded for predicting the majority
class. The regularization strength ($C = 0.1$ for both splits) and the threshold were chosen on
validation, then frozen.

\textbf{Graph neural network.} Twelve embedding tables of width 128 turn the category codes into
vectors. Three \texttt{GINEConv} layers \cite{hu2020,xu2019} pass messages between bonded atoms, each
applying batch normalization, a ReLU, dropout 0.3 and a residual connection;
\texttt{global\_mean\_pool} averages the atom vectors and a two-layer head produces one logit, for
242,564 trainable parameters. GINEConv adds the bond embedding to the neighbor's atom embedding
\emph{before} summing messages, so bond attributes modulate every message: an aromatic and a single bond
between the same elements send different messages. Three layers propagate information across a three-bond
radius, enough for a ring and its substituents, and mean pooling keeps a large molecule from producing a
large vector merely by having more atoms.

\textbf{Training.} The loss is binary cross-entropy weighted by \texttt{pos\_weight} =
negatives\,/\,positives from the training subset alone. AdamW runs at learning rate $10^{-3}$ with
cosine decay, weight decay $10^{-5}$ and batches of 128, for at most 150 epochs, stopping after 30
without a better validation ROC-AUC and restoring the best epoch's weights. Being threshold-free, ROC-AUC
keeps the choice of epoch independent of the later choice of threshold. Everything runs on the CPU, which
at this size beats the GPU (13\,s against 20\,s for 40 epochs) and, unlike GPU scatter-add, is exactly
reproducible. Validation ROC-AUC over three seeds is $0.920 \pm 0.007$ (random) and $0.911 \pm 0.004$
(scaffold).

\textbf{Reading the curves.} In Figure~\ref{fig:curves}, training loss falls smoothly and is still
falling when each run ends, while validation loss stops improving within about ten epochs and drifts
upward: ordinary overfitting, with no optimization instability. Validation loss does spike, by up to 0.2
between epochs, but that is a small-sample effect, since each validation subset holds under 300 molecules
and a few confidently wrong predictions move the mean. That is why the stopping rule watches validation
ROC-AUC, which plateaus near 0.90 well before the budget ends.

\begin{figure}[tb]
\centering\fig{E_curves}
\caption{\textbf{Both runs overfit early; validation ROC-AUC plateaus long before the epoch budget ends.}
Seed 0 of the random split (a, c) and the scaffold split (b, d). Everything here is the graph network,
so everything is red: training loss is the light tint, validation loss and ROC-AUC the full hue. Each row
shares a scale and both columns share an epoch axis, so the runs compare directly. Dashed line: the
restored epoch. Validation $n = 295$; loss is binary cross-entropy with the training-set positive
weight.}
\label{fig:curves}
\end{figure}

\section{Results}
\label{sec:results}

Each test subset was evaluated exactly once, at the threshold frozen on the matching validation subset.
Table~\ref{tab:comparison} reports every final evaluation; the network's entries are mean $\pm$ SD over
three seeds. Figures~\ref{fig:roc} and~\ref{fig:confusion} show the curves and the confusion matrices for
the median-validation seed, the one whose validation ROC-AUC is the middle of the three.

\begin{table}[t]
\centering\footnotesize\setlength{\tabcolsep}{3pt}
\caption{\textbf{Test results for both models under both splits.} Thresholds were fixed on the matching
validation subset by maximizing Youden's $J$ \cite{youden1950}, which is sensitivity plus specificity
minus one, and never adjusted afterwards. GNN rows are mean $\pm$ SD over
three seeds; fingerprint rows are a single deterministic fit. Acc.\ = accuracy, Sens.\ = sensitivity,
Spec.\ = specificity. $n = 294$ test molecules (random), 295 (scaffold).}
\label{tab:comparison}
\begin{tabular}{llccccccc}
\toprule
Model & Split & Thresh. & ROC-AUC $\uparrow$ & PR-AUC $\uparrow$ & Acc. $\uparrow$ & Sens. $\uparrow$ & Spec. $\uparrow$ & F1 $\uparrow$ \\
\midrule
Fingerprint & random & 0.496 & 0.919 & 0.970 & 0.867 & 0.893 & 0.783 & 0.912 \\
GNN & random & 0.574 $\pm$ 0.230 & 0.903 $\pm$ 0.009 & 0.959 $\pm$ 0.005 & 0.795 $\pm$ 0.021 & 0.772 $\pm$ 0.024 & 0.870 $\pm$ 0.014 & 0.852 $\pm$ 0.017 \\
Fingerprint & scaffold & 0.641 & 0.876 & 0.947 & 0.803 & 0.804 & 0.803 & 0.859 \\
GNN & scaffold & 0.843 $\pm$ 0.227 & 0.866 $\pm$ 0.011 & 0.947 $\pm$ 0.006 & 0.783 $\pm$ 0.051 & 0.781 $\pm$ 0.086 & 0.789 $\pm$ 0.053 & 0.841 $\pm$ 0.047 \\
\bottomrule
\end{tabular}

\end{table}

\textbf{The graph network does not beat the fingerprint baseline.} The baseline scores higher under both
splits, by 0.016 ROC-AUC under the random split (0.919 against $0.903 \pm 0.009$) and by 0.011 under the
scaffold split (0.876 against $0.866 \pm 0.011$). Neither difference is statistically distinguishable.
The baseline's own 95\% bootstrap intervals, $[0.876, 0.955]$ and $[0.824, 0.922]$, are roughly
$\pm 0.04$, and the network's three seeds alone differ by as much as the scaffold-split gap. Those
intervals are \emph{marginal}: each carries the variation from this test subset being easy or hard, which
is shared between the models and cancels out of a difference. Scoring both on the same resampled
molecules is the stricter test and puts the difference at $-0.025$ $[-0.056, +0.003]$ (random) and
$-0.016$ $[-0.057, +0.026]$ (scaffold), both containing zero; these use the median-validation seed, not
the three-seed mean. The evidence supports no ranking between the models, worth stating because the
opposite is commonly assumed. A graph network reads connectivity, which a fingerprint discards, but with
only 1,375 training molecules it has too little data to use that extra expressiveness, while a 2,048-bit
fragment list already captures most of what separates these classes.

\textbf{Changing the split moves the score more than changing the model does.} Both models score lower
under the scaffold split, the network by 0.037 and the baseline by 0.043, against a seed-to-seed spread
of at most 0.011. The split effect is three to four times the noise; the model effect is not. Section~4
gives the mechanism, and the simpler model is flattered slightly more by the random split, as expected if
those points come from recognizing familiar families rather than chemistry.

\textbf{Reading the operating points.} The network's thresholds are unstable ($0.574 \pm 0.230$ and
$0.843 \pm 0.227$) even though the resulting sensitivity and specificity are not, because the ROC curve
is flat near its optimum and many thresholds earn almost the same Youden's $J$. At its chosen point the
scaffold-split network finds 78.1\% of penetrating compounds while rejecting 78.9\% of non-penetrating
ones; the baseline is at 80.4\% and 80.3\%. PR-AUC is high everywhere (0.947--0.970) because three
quarters of the molecules are positive, which puts its chance line at the test prevalence, 0.765 (random)
and 0.742 (scaffold), not 0.5. It is not the headline.

\begin{figure}[tb]
\centering\fig{F_curves}
\caption{\textbf{All four configurations separate the classes well, and the curves are hard to tell
apart.} Test (a) ROC and (b) precision--recall curves. Color encodes the model (red = GNN, petrol =
fingerprint), lightness the split (full = scaffold, pale = random), line style the model again so pairs
stay distinct in grayscale, which is what separates them there rather than lightness. The key gives each
curve's ROC-AUC and PR-AUC, with $\pm$ only on the network, whose three seeds differ; the baseline is one
deterministic fit. Those numbers matter because the precision--recall curves cross repeatedly between
recall 0.2 and 0.6, where which curve is in front says nothing about which model scores higher. The
precision axis is clipped to the range the curves occupy, and the dashed chance line in (b) sits at the
mean test prevalence (0.765 random, 0.742 scaffold). $n = 294$, $n = 295$.}
\label{fig:roc}
\end{figure}

\begin{figure}[tb]
\centering\fig{F_confusion}
\caption{\textbf{At their chosen thresholds all four configurations trade sensitivity for specificity in
much the same way.} Test confusion matrices, counts with row percentages, each at that configuration's
own validation-fixed threshold (Table~\ref{tab:comparison}); GNN panels show the median-validation seed.
Hue names the model. Shading encodes the row share alone, on one scale across all four panels, so a
cell's depth compares between them.}
\label{fig:confusion}
\end{figure}

\section{Where the model goes wrong}

The errors examined are the network's on the scaffold split, at the threshold of 0.967 that this run's
validation subset selected; Table~\ref{tab:comparison} gives $0.843 \pm 0.227$ because that is the
three-seed mean, and the gap between them is the threshold instability noted in
Section~\ref{sec:results}. That threshold leaves 16 false positives and 42 false negatives among 295
molecules. The four in Figure~\ref{fig:errors} follow a fixed rule: the two false positives and two false
negatives the model was most confident about, each paired with its nearest training molecule by Tanimoto
similarity (Table~\ref{tab:errors}).

\textbf{Similarity explains one of the four errors, and the other three in the opposite way.} The first
false positive has a training neighbor at Tanimoto 0.753 carrying the \emph{opposite} label: two
molecules sharing almost every fragment, labeled differently. A model reading fragments cannot separate
them, so that error belongs to the data rather than the network. The other three are the reverse case,
with nearest neighbors at Tanimoto 0.273, 0.248 and 0.172. Those molecules sit outside the chemistry the
model was trained on, so the predictions are extrapolations, exactly what the scaffold split was built to
create. The model was \emph{confident} in all four ($p = 1.000$, 1.000, 0.006, 0.008) and gives no signal
that it is extrapolating, which matters more for deployment than the error rate. A threshold of 0.967
points the same way: probabilities piled against 0 and 1 are poorly calibrated, so these are ranking
scores rather than chances of penetration.

\begin{figure}[tb]
\centering\fig{G_errors}
\caption{\textbf{Three of the four most confident errors are extrapolations away from the training
chemistry; one is a near-duplicate with the opposite label.} The network's two most confident false
positives (top) and false negatives (bottom) on the scaffold-split test set, each paired to its right with
its nearest training molecule by Tanimoto similarity; a hairline divides the left pair from the right and
the labels FP~1, FP~2, FN~1, FN~2 match Table~\ref{tab:errors}. Each tag gives the true class and, for
test molecules, the predicted probability beneath it; training molecules are prefixed ``train''.
Structures share one bond length, so sizes are comparable. Threshold 0.967, the median-validation seed's
own value.}
\label{fig:errors}
\end{figure}

\section{Limitations}
\label{sec:limits}

\textbf{The label is uncertain, and this dataset shows it.} \texttt{p\_np} records what curated
literature concluded, from studies differing in method, species and dose; ten canonical molecules carry
both labels in different records (Section~\ref{sec:clean}). Those are only the disagreements visible
\emph{inside} one file; nothing measures how often the surviving label is wrong, so
a model fitted to it can be no more reliable, and ROC-AUC against that label cannot reveal the
ceiling.

\textbf{The dataset is small and chemically narrow.} It holds 1,965 molecules, 76\% of them positive,
and 777 of its scaffolds appear once. The bootstrap intervals are correspondingly wide, about $\pm 0.04$
ROC-AUC --- precisely the size of the model differences the experiment was built to test. The compounds
come from published permeability studies, which over-represent marketed drugs and their analogs.

\textbf{Nothing was validated outside this dataset.} Every number comes from one curated file split two
ways: nothing was measured in a laboratory, no external dataset served as a check, no prediction was
tested prospectively. The scaffold split is the strongest available substitute, but still a rearrangement
of the same 1,965 molecules.

\looseness=-1
\textbf{Why this establishes neither causality nor clinical suitability.} The model learns which atom
and bond patterns co-occur with the positive label. Co-occurrence is not mechanism: barrier penetration
involves passive permeability, active transport and efflux, none of which appears in a graph of atoms and
bonds. A high score says a molecule resembles the ones labeled penetrating here, not why and not what it
would do in a person. The proper use is to order candidates for laboratory work, leaving assays, not
model scores, as the basis for decisions.

\section{Conclusion}

On BBBP the graph network does not outperform the fingerprint baseline, and a paired bootstrap cannot
separate them under either split. The splitting strategy moves
the score three to four times as much as the model does, so any BBBP number quoted without it is
uninterpretable. Two things follow. At this dataset size a fingerprint baseline is the model to beat, not
a formality. And the network is most confident, at $p = 1.000$ and $0.006$, exactly where its nearest
training neighbor shares almost no structure: a tool that cannot signal when it is extrapolating is
harder to deploy safely than one scoring lower, so calibration comes first.

\bibliographystyle{unsrt}
\begin{thebibliography}{9}
\bibitem{wu2018} Z.~Wu, B.~Ramsundar, E.~N.~Feinberg, J.~Gomes, C.~Geniesse, A.~S.~Pappu, K.~Leswing and
V.~Pande. MoleculeNet: a benchmark for molecular machine learning. \emph{Chemical Science}, 9(2):513--530,
2018. doi:10.1039/C7SC02664A.
\bibitem{martins2012} I.~F.~Martins, A.~L.~Teixeira, L.~Pinheiro and A.~O.~Falcão. A Bayesian Approach to
in Silico Blood-Brain Barrier Penetration Modeling. \emph{Journal of Chemical Information and
Modeling}, 52(6):1686--1697, 2012. doi:10.1021/ci300124c.
\bibitem{fey2019} M.~Fey and J.~E.~Lenssen. Fast graph representation learning with PyTorch Geometric. In
\emph{ICLR Workshop on Representation Learning on Graphs and Manifolds}, 2019. arXiv:1903.02428.
\bibitem{bemis1996} G.~W.~Bemis and M.~A.~Murcko. The Properties of Known Drugs. 1. Molecular Frameworks.
\emph{Journal of Medicinal Chemistry}, 39(15):2887--2893, 1996. doi:10.1021/jm9602928.
\bibitem{rogers2010} D.~Rogers and M.~Hahn. Extended-Connectivity Fingerprints. \emph{Journal of Chemical
Information and Modeling}, 50(5):742--754, 2010. doi:10.1021/ci100050t.
\bibitem{hu2020} W.~Hu, B.~Liu, J.~Gomes, M.~Zitnik, P.~Liang, V.~Pande and J.~Leskovec. Strategies for
Pre-training Graph Neural Networks. In \emph{International Conference on Learning Representations}, 2020.
arXiv:1905.12265.
\bibitem{xu2019} K.~Xu, W.~Hu, J.~Leskovec and S.~Jegelka. How powerful are graph neural networks? In
\emph{International Conference on Learning Representations}, 2019. arXiv:1810.00826.
\bibitem{rdkit} RDKit: Open-source cheminformatics, version 2025.03.6. \url{https://www.rdkit.org},
doi:10.5281/zenodo.591637. (The project states it has no citable publication and gives this form.)
\bibitem{youden1950} W.~J.~Youden. Index for rating diagnostic tests. \emph{Cancer}, 3(1):32--35, 1950.
doi:10.1002/1097-0142(1950)3:1<32::AID-CNCR2820030106>3.0.CO;2-3.
\end{thebibliography}

\clearpage
\appendix
% Appendix tables are numbered A1, A2, A3 so that the one table in the main text is simply Table 1 and a
% reader is never sent from Section 3 to a "Table 3" that appears six pages later.
\setcounter{table}{0}
\renewcommand{\thetable}{A\arabic{table}}
\section*{Appendix}

\begin{table}[h]
\centering\footnotesize
\caption{Split composition. The random split is stratified, so its positive proportion is constant; the
scaffold split has no mechanism to control the balance yet stays close to it. Scaffold counts apply only
to the scaffold split, whose three scaffold sets are disjoint.}
\label{tab:splits}
\begin{tabular}{llccc}
\toprule
Split & Subset & Molecules & Positive & Unique scaffolds \\
\midrule
random & train & 1376 & 1050 (0.763) & -- \\
random & validation & 295 & 225 (0.763) & -- \\
random & test & 294 & 225 (0.765) & -- \\
scaffold & train & 1375 & 1054 (0.767) & 430 \\
scaffold & validation & 295 & 227 (0.769) & 295 \\
scaffold & test & 295 & 219 (0.742) & 295 \\
\bottomrule
\end{tabular}

\end{table}

\begin{table}[h]
\centering\footnotesize
\caption{The twelve categorical feature columns of the PyTorch Geometric graphs. ``Vocabulary'' is the
number of categories the encoding allows; ``codes seen'' is how many actually occur in BBBP. Each column
is encoded by its own trainable embedding table.}
\label{tab:features}
\begin{tabular}{llcc}
\toprule
Column & Meaning & Vocabulary & Codes seen \\
\midrule
\multicolumn{4}{l}{\itshape Atom features, \texttt{x}} \\
\texttt{atomic\_num} & chemical element (code = atomic number) & 119 & 13 \\
\texttt{chirality} & local three-dimensional arrangement at the atom & 9 & 3 \\
\texttt{degree} & number of bonded neighbors in the graph & 11 & 5 \\
\texttt{formal\_charge} & assigned charge (code 5 is neutral) & 12 & 4 \\
\texttt{num\_hs} & attached hydrogens, which the graph does not store as nodes & 9 & 4 \\
\texttt{num\_radical\_electrons} & unpaired electrons & 5 & 2 \\
\texttt{hybridization} & orbital hybridization (sp, sp2, sp3, ...) & 8 & 4 \\
\texttt{is\_aromatic} & atom belongs to an aromatic ring & 2 & 2 \\
\texttt{is\_in\_ring} & atom belongs to any ring & 2 & 2 \\
\addlinespace
\multicolumn{4}{l}{\itshape Bond features, \texttt{edge\_attr}} \\
\texttt{bond\_type} & single, double, triple or aromatic bond & 22 & 4 \\
\texttt{stereo} & double-bond stereochemistry (none, Z, E, ...) & 6 & 3 \\
\texttt{is\_conjugated} & bond is part of a conjugated system & 2 & 2 \\
\bottomrule
\end{tabular}

\end{table}

\begin{table}[h]
\centering\footnotesize
\caption{The four most confident errors of the graph network on the scaffold-split test set, with the
nearest training molecule by Morgan-fingerprint Tanimoto similarity. ``True'' and ``nearest train label''
are BBBP labels (1 = penetrating). Full SMILES are in the notebook.}
\label{tab:errors}
\begin{tabular}{llcccc}
\toprule
Case & Error & True & Predicted $p$ & Nearest train label & Tanimoto \\
\midrule
FP 1 & false positive & 0 & 1.000 & 1 & 0.753 \\
FP 2 & false positive & 0 & 1.000 & 0 & 0.273 \\
FN 1 & false negative & 1 & 0.006 & 0 & 0.248 \\
FN 2 & false negative & 1 & 0.008 & 1 & 0.172 \\
\bottomrule
\end{tabular}

\end{table}

\vspace{1em}
\noindent\end{document}
